% !TeX root = main.tex
\chapter{Модель мультимодального представлення прихованої загрози в мультимедійних об’єктах}
\label{chap:models_formalization}

Проведений у першому розділі аналіз показав, що прихований вплив у мультимедійному об’єкті може залишати ознаки в декодованому вмісті, частотному або хвильковому представленні, структурі контейнера, службових даних і перебігу контрольованого оброблення. Окреме відхилення, однак, не доводить наявності аномального коду: воно може бути наслідком допустимого перетворення, особливості формату або неповного спостереження.

Звідси постає дослідницька суперечність: виявлення слабких прихованих змін вимагає спільного врахування різнорідних ознак, тоді як їх об’єднання не повинно посилювати висновок понад фактично доступні свідчення. Просте накопичення показників цієї суперечності не усуває, оскільки приховує походження ознаки, умови її одержання та причину відсутності окремого спостереження. Для розв’язання суперечності розроблено модель мультимодального представлення прихованої загрози, у якій усі складники підпорядковано єдиному правилу збереження походження свідчень.

Побудова моделі відтворює послідовність дослідження. Спочатку визначено умови, за яких мультимедійний об’єкт проходить штатний канал приймання. Далі встановлено, які представлення об’єкта необхідно зберегти, щоб не втратити початкові байти, декодований вміст, будову контейнера та причинно пов’язані події. Після цього різнорідні спостереження перетворюються на ознаки зі збереженням координат і походження, а на завершальному кроці узгоджене представлення переходить в аналітичний стан, що містить оцінку ризику, упевненість і пояснення. Перелічені кроки є взаємозалежними частинами однієї моделі, а не переліком окремих моделей.

У роботі розмежовано чотири базові поняття, кожне з яких має власний обсяг. \textit{Аномалією} названо спостережуване відхилення від установленої структурної, статистичної або поведінкової норми. \textit{Аномальним кодом} є цілеспрямовано сформована послідовність даних, команда, фрагмент коду або структурне значення, що не відповідає заявленому призначенню носія та може підтримати несанкціоновану дію. \textit{Прихована загроза} описує можливий небезпечний вплив, для якого встановлюється зв’язок між властивостями носія, операцією оброблення та очікуваним наслідком. Термін \textit{шкідливий код} застосовується лише тоді, коли підтверджено програмний зміст і шкідливу функцію. Додатково розрізнено два супровідні поняття: \textit{приховане навантаження} позначає перенесені носієм дані, а \textit{прихований вплив} --- можливий або фактичний наслідок їх опрацювання компонентом веб-ресурсу.

Модель спирається на п’ять припущень: початковий мультимедійний об’єкт зберігається незмінним; для кожного похідного представлення відоме правило його одержання; відсутність представлення не ототожнюється з нульовим ризиком; подія належить об’єкту лише за підтвердженого причинного зв’язку; одержаний висновок характеризує стан за доступних даних і чинних налаштувань, але не гарантує абсолютної безпечності. Саме ці припущення визначають силу подальших тверджень і межі застосування моделі.

Межу довіри проведено між фактом збереження початкової байтової послідовності та результатами її тлумачення. Послідовність $X_{mm}$ достовірно відтворює те, що надійшло до перевірки, але сама по собі не підтверджує ні заявленого формату, ні безпечності вмісту. Засіб структурного аналізу, засіб декодування та контрольоване середовище мають власні обмеження і можуть завершити роботу неповним результатом, тому разом з ознаками зберігаються версії засобів, показники якості, правило походження даних і причина неповноти.

Два однакові числові результати можна порівнювати лише за сумісних схем ознак і умов спостереження. Зміна засобу декодування, правила приймання, маршруту або доступності журналу здатна змінити тлумачення ознаки без зміни початкового об’єкта. Повторна перевірка через це доповнює історію спостережень, а не заміщує попередній результат без збереження його умов.

Щоб не переобтяжувати виклад, математичні залежності подано для одного мультимедійного об’єкта без наскрізного індексу. Під час опису вибірки або повторних запусків індекс $i$ додається до відповідного позначення, наприклад $X_{mm,i}$ або $Z_i$, без зміни його змісту.

\section{Умови здійснення атак на веб-ресурси з використанням мультимедійних об’єктів}
\label{sec:attack_model}

Мультимедійний об’єкт у досліджуваному сценарії є не кінцевою ціллю атаки, а носієм прихованого навантаження та засобом його передавання через штатний або формально дозволений канал. Небезпека виникає з різниці правил тлумачення: об’єкт може задовольняти перевірку під час приймання, але залишатися неоднозначним для наступних компонентів, які опрацьовують декодовані значення, структуру контейнера, службові поля або часову організацію відео.

Наперед перелічити всі способи внесення прихованих даних неможливо, а прив’язування моделі до одного способу зробило б її непридатною до нових варіантів підготовки носія. Тому в моделі відокремлено дію порушника від наслідків, які спостерігаються під час перевірки. Узагальнений зв’язок між початковим об’єктом, прихованим навантаженням і підготовленим носієм подано формулою~\eqref{eq:attack_embedding}:
\begin{equation}
X_{mm}^{\prime}=E_A(X_{mm},P_A,\Theta_A),
\label{eq:attack_embedding}
\end{equation}
де, $X_{mm}$ --- початковий мультимедійний об’єкт; $P_A$ --- приховане навантаження; $\Theta_A$ --- параметри підготовки; $E_A$ --- функція вбудовування або цілеспрямованої зміни; $X_{mm}^{\prime}$ --- підготовлений мультимедійний об’єкт.

Формула~\eqref{eq:attack_embedding} фіксує межу дослідження: у ній описано результат цілеспрямованої підготовки, але не закладено конкретної техніки внесення даних. Завдяки цьому подальший аналіз ґрунтується на спостережуваних властивостях $X_{mm}^{\prime}$, а не на припущенні про заздалегідь відомий спосіб атаки. Залежність не є окремим методом дисертації й не означає, що зовнішній вигляд об’єкта обов’язково змінюється.

Проходження підготовленого носія до точки спостереження розглянуто як послідовність чотирьох умов.

\begin{enumerate}
\item \textit{Проходження правила приймання.} Значення $V_f(X_{mm}^{\prime})=1$, де $V_f$ --- бінарна функція чинної перевірки формату, означає, що об’єкт визнано припустимим за правилом приймання. Виконання цієї умови лише відкриває подальший маршрут і не підтверджує безпечності вмісту.

\item \textit{Збереження прихованості та навантаження.} Зовнішнє відхилення обмежується умовою $D(X_{mm},X_{mm}^{\prime})\leq\varepsilon$, а збереження значущої частини навантаження на маршруті $G$, означеному в четвертій умові, --- нерівністю $S(P_A,G)\geq\tau_A$. Тут $D$ --- функція відхилення між об’єктами; $\varepsilon$ --- допустимий рівень змін; $S$ --- оцінка збереження навантаження; $\tau_A$ --- її поріг. Виконання цих умов характеризує можливість доставки, але не доводить прояву впливу.

\item \textit{Проходження компонентів веб-ресурсу.} Функціонально відмінні компоненти утворюють скінченну множину $R_{web}=\{r_1,r_2,\ldots,r_{n_R}\}$, до якої можуть належати засоби приймання, первинної перевірки, структурного розбору, декодування, перетворення, збереження, відображення та контрольованого спостереження. Для дослідження істотне не фізичне розміщення компонента, а правило, за яким він тлумачить об’єкт.

\item \textit{Відтворення фактичного маршруту.} Порядок проходження задається послідовністю $G=\langle g_1,g_2,\ldots,g_{n_G}\rangle$, кожний елемент якої $g_j\in R_{web}$ відповідає зафіксованому кроку маршруту; той самий компонент може траплятися в маршруті неодноразово. Разом із кроком зберігаються стан компонента $\theta_j$, дозволена операція та час спостереження. Це дає змогу віднести подію до певної фази оброблення, а не лише до спільного часового проміжку.
\end{enumerate}

Наведена послідовність відокремлює три різні факти: наявність прихованого навантаження, його збереження під час проходження та спостережуваний прояв впливу. Структурне відхилення характеризує властивість контейнера; проходження до певного компонента свідчить про збереження цієї властивості на відповідній ділянці маршруту; лише причинно пов’язана реакція дає підставу говорити про можливий прояв. Жоден із цих фактів окремо не встановлює наміру порушника або програмного змісту фрагмента.

Якщо відомості про частину маршруту відсутні, модель не доповнює їх припущеними кроками: невідома ділянка фіксується як межа причинного висновку. Такий підхід зменшує ризик приписати мультимедійному об’єкту сторонню подію лише через часовий збіг.

Причинну послідовність від підготовки носія до одержання доступних спостережень узагальнено на рис.~\ref{fig:attack_model}. Схема показує, чому під час аналізу слід зберігати не лише ознаку, а й точку маршруту, у якій вона виникла.

\begin{figure}[htbp]
\centering
\begin{tikzpicture}[x=1cm,y=1cm]
\node[dissertation block,text width=2.75cm] (object) at (-5.65,1.35)
  {Початковий об’єкт\\$X_{mm}$};
\node[dissertation block,text width=2.75cm] (payload) at (-5.65,-0.65)
  {Приховане навантаження\\$P_A$};
\node[dissertation block,text width=3.00cm] (prepare) at (-2.10,0.35)
  {Підготовка об’єкта\\$E_A(\Theta_A)$};
\node[dissertation block,text width=2.75cm] (prepared) at (1.20,0.35)
  {Підготовлений об’єкт\\$X_{mm}^{\prime}$};
\node[dissertation block,text width=3.15cm] (conditions) at (4.65,0.35)
  {Умови прихованості\\$V_f=1$, $D\leq\varepsilon$};
\node[dissertation block,text width=3.30cm] (route) at (4.65,-1.75)
  {Маршрут $G$\\$r_j\in R_{web}$};
\node[dissertation block,text width=4.40cm] (observations) at (0.55,-3.20)
  {Статичні та подієві спостереження\\$X_s,X_v,X_f,X_b$};

\coordinate (inputbus) at (-4.05,0.35);
\draw[thick] (object.east) -- (object.east -| inputbus) -- (inputbus);
\draw[thick] (payload.east) -- (payload.east -| inputbus) -- (inputbus);
\fill[black] (inputbus) circle (1.3pt);
\draw[dissertation arrow] (inputbus) -- (prepare.west);
\draw[dissertation arrow] (prepare.east) -- (prepared.west);
\draw[dissertation arrow] (prepared.east) -- (conditions.west);
\draw[dissertation arrow] (conditions.south) -- (route.north);
\draw[dissertation arrow] (route.south) |- (observations.east);
\end{tikzpicture}
\caption{Причинна схема формування спостережень за умов прихованої доставки мультимедійного об’єкта}
\label{fig:attack_model}
\end{figure}


Установлені умови описують можливість прихованої доставки, але не підміняють ознакового аналізу твердженням про виконання шкідливого коду. Для перевірки потрібні представлення, які одночасно зберігають початковий носій, результати його контрольованого тлумачення та межі доступного спостереження.

\section{Формалізація мультимедійного об’єкта як багаторівневого інформаційного контейнера}
\label{sec:multilevel_container}

Один мультимедійний об’єкт опрацьовується за різними правилами, і кожне похідне представлення втрачає частину початкової інформації. Засіб декодування робить вміст придатним для зорового й часового аналізу, але може не відображати службові ділянки або самочинно виправляти структурні помилки. Побайтовий розбір зберігає будову контейнера, проте не пояснює його зорового змісту. Журнал подій показує реакцію середовища, але не доводить, що кожна подія спричинена саме досліджуваним об’єктом. Через це одиницею аналізу обрано одну простежувану перевірку, яка поєднує представлення без їх взаємної підміни.

Ідентифікатор перевірки пов’язує початковий об’єкт, похідні дані, ознаки, події та результат. Тотожність входу підтверджується незмінністю байтової послідовності, а належність похідних даних --- збереженим правилом їх одержання. Це переносить межу довіри до початкових байтів і дає змогу розглядати решту даних як спостереження, для яких зазначено використані засоби та їх версії.

Формування повного входу моделі виконується у п’ять взаємозалежних кроків.

\begin{enumerate}
\item \textit{Збереження початкового носія.} Об’єкт $X_{mm}=(x_1,x_2,\ldots,x_n)$ зберігається як послідовність $n$ байтів зі значеннями з множини $\{0,\ldots,255\}$. Очищення, перекодування або автоматичне виправлення створює похідну копію і не заміщує $X_{mm}$. Це правило принципове, оскільки саме перетворення може знищити досліджуване відхилення.

\item \textit{Одержання зорових і часових даних.} Контрольоване декодування утворює $X_v=E_v(X_{mm},\theta_v)$, де $E_v$ --- засіб перетворення, а $\theta_v$ містить його параметри та версію. Для зображення $X_v$ містить один кадр, для відео --- впорядковану послідовність кадрів. Частотне представлення $X_f$ і часове $X_t$ створюються пізніше, під час виділення ознак, тому вони не замінюють декодованого входу.

\item \textit{Відтворення будови контейнера.} Структурне представлення $X_s=\langle s_1,s_2,\ldots,s_{n_s}\rangle$ зберігає порядок і вкладеність установлених елементів. Запис $s_k=\langle c_k,a_k,b_k,w_k,v_k\rangle$ містить тип елемента $c_k$, початкове та кінцеве байтові зміщення $a_k$ і $b_k$, посилання $w_k$ на батьківський елемент і результат $v_k$ місцевої перевірки форматних обмежень. Пов’язані службові значення утворюють $X_m$. Завдяки координатам будь-яке структурне відхилення можна повернути до ділянки початкового носія.

\item \textit{Відбір причинно пов’язаних подій.} Подієво-поведінкове представлення $X_b=\langle e_1,e_2,\ldots,e_T\rangle$ містить лише події, належність яких до поточної перевірки підтверджено. Для події $e_t$ зберігаються дія, ініціатор, ресурс, результат і час. Фаза оброблення, ідентифікатор запуску та стан середовища додаються до супровідного запису. Сам часовий збіг не є достатньою підставою для включення події до $X_b$.

\item \textit{Фіксація доступності та якості.} Для трьох гілок зберігається маска доступності $m_A=(m_s,m_v,m_b)\in\{0,1\}^3$, складники якої відповідають структурним, спектрально-зоровим і подієво-поведінковим даним. Нульове значення означає відсутність придатного спостереження, а не підтвердження норми. Причина недоступності, повнота спостереження та версії засобів входять до контексту $C$.
\end{enumerate}

Окремі представлення мають цінність лише тоді, коли їх можна віднести до одного входу та відтворити умови одержання. Статичні дані попередньо згруповано як $X_c=\langle X_{mm},X_v,X_s,X_m\rangle$. Узгоджений склад повного входу подальшого аналізу визначено формулою~\eqref{eq:prepared_packet}:
\begin{equation}
X_A=\langle X_c,X_b,m_A,C\rangle,
\label{eq:prepared_packet}
\end{equation}
де, $X_A$ --- повний вхід моделі; $X_c$ --- сукупність статичних даних; $X_b$ --- причинно пов’язані події; $m_A$ --- маска доступності; $C$ --- відомості про походження, параметри та якість.

Формула~\eqref{eq:prepared_packet} пов’язує всі спостереження з однією перевіркою і водночас залишає їх розділеними за природою. Такий поділ обумовлює два подальші рішення: статичні ознаки формуються без домішування реакції середовища, а поведінкові приєднуються лише після підтвердження їх причинної належності. Склад даних не задає конкретної програмної реалізації, а встановлює вимоги до відтворюваності аналізу.

Граничні випадки фіксуються явно. Якщо причинно пов’язані події недоступні, $X_b$ не містить придатних записів, а $m_b=0$. Якщо декодування безпечно припинилося, структурна гілка може залишатися доступною за $m_v=0$. Якщо походження похідного представлення не підтверджено, його не включають до $X_A$ навіть за технічної можливості прочитати значення. Доступне, але неповне представлення супроводжується показником покриття й обмеженням тлумачення.

Походження статичних і подієво-поведінкових складників показано на рис.~\ref{fig:multilevel_container}. Схема пояснює, через що початковий об’єкт, похідні представлення та події не можна передавати до ознакового аналізу без відомостей про доступність і походження.

\begin{figure}[htbp]
\centering
\begin{tikzpicture}[x=1cm,y=1cm]
\node[dissertation block,text width=3.05cm] (object) at (-5.70,1.20)
  {Незмінний об’єкт\\$X_{mm}$};
\node[dissertation block,text width=3.35cm] (visual) at (-2.05,2.25)
  {Зорові дані\\$X_v$};
\node[dissertation block,text width=3.35cm] (structural) at (-2.05,0.15)
  {Структурні дані\\$X_s,X_m$};
\node[dissertation block,text width=4.15cm] (static) at (1.70,1.20)
  {Статичні дані після попереднього опрацювання\\$X_c=\langle X_{mm},X_v,X_s,X_m\rangle$};
\node[dissertation block,text width=3.40cm] (events) at (-2.05,-2.20)
  {Причинно пов’язані події\\$X_b$};
\node[dissertation block,text width=3.20cm] (context) at (1.70,-1.35)
  {Доступність і супровідні\\відомості $(m_A,C)$};
\node[dissertation block,text width=3.20cm] (input) at (5.45,0.15)
  {Повний вхід моделі\\$X_A$};

\draw[dissertation arrow] (object.east) -- (visual.west);
\draw[dissertation arrow] (object.east) -- (structural.west);
\draw[dissertation arrow] (object.north) -- ++(0,1.70) -| (static.north);
\draw[dissertation arrow] (visual.east) -- (static.west);
\draw[dissertation arrow] (structural.east) -- (static.west);
\draw[dissertation arrow] (static.east) -- (input.west);
\draw[dissertation arrow] (events.east) -- ++(0.55,0) -| ([xshift=-0.45cm]input.south);
\draw[dissertation arrow] (context.east) -- (input.west);
\end{tikzpicture}
\caption{Походження статичних і подієво-поведінкових даних повного входу моделі}
\label{fig:multilevel_container}
\end{figure}


Таким чином, багаторівневе представлення узгоджує повноту початкового носія з вибірковістю похідних спостережень: зберігаються незмінний вхід, координати відхилень, причинна належність подій і явна неповнота. Наступний крок полягає в перетворенні цих даних на порівнювані ознаки без втрати їхнього походження.

\section{Ознакове ядро моделі мультимодального представлення}
\label{sec:feature_space}

Багаторівневі дані мають різну фізичну природу й різні системи координат. Їх передчасне об’єднання створює ризик числового переважання однієї групи ознак і розриває зв’язок між узагальненою оцінкою та спостережуваною ділянкою. Через це ознакове ядро має не лише привести дані до порівнюваного вигляду, а й зберегти їх відмінність до моменту обґрунтованого узгодження.

Перетворення даних на ознаки організовано за чотирма послідовними кроками.

\begin{enumerate}
\item \textit{Формування структурних ознак.} Карта $X_s$ має змінну довжину й залежить від формату, тому безпосереднє порівняння об’єктів за нею неможливе. Функція $F_s$ перетворює $X_s$ і службові значення $X_m$ на числове представлення $Z_s=F_s(X_s,X_m)$. Паралельно функція $G_s$ утворює множину байтових координат $L_s=G_s(X_s)$. Числові ознаки потрібні для зіставлення, а координати --- для перевірки, які ділянки початкового носія підтримують одержану оцінку.

\item \textit{Формування спектрально-зорових ознак.} Структурно коректний контейнер може містити нетипові закономірності в декодованому вмісті. Тому зорове представлення доповнюється частотним $X_f$ і, для відео, часовим $X_t$ представленням. Функція $F_v$ формує $Z_v=F_v(X_v,X_f,X_t)$, а функція $G_v$ зберігає просторово-часові координати $L_v=G_v(X_v,X_t,\tau_v)$, де $\tau_v$ --- локальний поріг відбору локалізацій, який належить до параметрів першого методу й конкретизується в розділі~\ref{chap:methods_technology}. Для зображення відсутність часової складової є властивістю об’єкта, а не технічною відмовою.

\item \textit{Формування подієво-поведінкових ознак.} Послідовність $X_b$ має змінну довжину й містить категорійні та часові значення. Після перевірки причинної належності подій функція $F_b$ формує $Z_b=F_b(X_b)$. Відсутність події має зміст лише тоді, коли відповідна фаза була доступною для спостереження. Тому $Z_b$ завжди розглядається разом із $m_b$ і показниками повноти журналу.

\item \textit{Узгодження походження.} Байтові координати $L_s$ і просторово-часові координати $L_v$ зіставляються лише через збережене правило одержання декодованого представлення. Відсутність прямої відповідності не спростовує жодної гілки, оскільки декодування може мати нелокальний характер. Вона лише обмежує силу твердження про спільне джерело відхилень.
\end{enumerate}

Для обчислення спільної статичної оцінки може використовуватися похідне представлення $Z_c=F_C(Z_v,Z_s)$; саме воно застосовується правиловою конфігурацією другого методу в підрозділі~\ref{sec:behavior_alignment}. Похідне представлення не замінює окремих $Z_v$ і $Z_s$, оскільки однакове значення $Z_c$ може бути одержане за різних поєднань ознак і різної повноти спостережень. Тому на межі між моделлю та методами передається статичний опис, склад якого наведено у формулі~\eqref{eq:static_profile}:
\begin{equation}
H=\langle Z_v,Z_s,L_v,L_s,m_c,C\rangle,
\label{eq:static_profile}
\end{equation}
де, $H$ --- статичний опис об’єкта; $Z_v$ і $Z_s$ --- спектрально-зорові та структурні ознаки; $L_v$ і $L_s$ --- їх просторово-часові та байтові координати; $m_c=(m_s,m_v)$ --- маска доступності статичних гілок; $C$ --- відомості про умови формування.

Статичний опис зберігає пояснюваність під час переходу до другого методу: він передає не лише числові значення, а й місця їх походження та ознаку доступності. Представлення $X_b$ до $H$ не входить, оскільки передчасне додавання подій змішало б властивості носія з реакцією середовища.

Після окремого формування трьох груп постає центральне питання моделі: як об’єднати їх, не прирівнюючи відсутності до нуля, а подібності --- до причинності. Відповіддю є узгодження лише тих ознак, що належать одному входу, одержані за сумісних умов і мають простежуване походження. Зв’язок між статичним описом, подієво-поведінковими ознаками та повним представленням наведено у формулі~\eqref{eq:integrated_representation}:
\begin{equation}
Z=F_I(H,Z_b,m_A,C)=\langle z,L,m_A,C\rangle,
\label{eq:integrated_representation}
\end{equation}
де, $Z$ --- інтегроване представлення; $F_I$ --- функція узгодженого об’єднання; $z$ --- спільний числовий вектор, який зберігає іменовані складники окремих гілок, зокрема їхні часткові оцінки, без злиття в одне значення; $L$ --- сукупність байтових і просторово-часових координат із посиланнями на події; $m_A$ --- маска доступності; $C$ --- відомості про якість і походження.

Формула~\eqref{eq:integrated_representation} задає не механічне усереднення, а правило збереження підстав результату: узгоджене підтвердження, взаємодоповнення, суперечність і недостатність спільних підстав залишаються різними станами. Якщо дві перевірки мають однаковий числовий вектор $z$, але в одній його підтримано кількома узгодженими гілками, а в іншій він виник за неповних або суперечливих даних, їхні висновки не вважаються рівносильними; різницю зберігають $L$, $m_A$ і $C$.

Об’єднання не приписує причинності лише на основі подібності ознак. Для причинного тлумачення потрібні спільний початковий об’єкт, сумісний маршрут і підтверджена послідовність подій. Якщо хоча б одну умову не виконано, ознака може залишатися описовою, але не посилює твердження про прояв прихованого впливу.

Послідовність формування $Z$ подано на рис.~\ref{fig:feature_space}. Схема пояснює, чому втрата однієї гілки позначається маскою доступності, а не приховано компенсується іншими ознаками.

\begin{figure}[H]
\centering
\begin{tikzpicture}[x=1cm,y=1cm,scale=0.88,transform shape,
  decision/.style={diamond,draw=black,line width=0.8pt,fill=white,
    aspect=1.35,align=center,font=\footnotesize,inner sep=1pt}]
\node[dissertation small,text width=3.10cm] (static) at (-6.50,2.20)
  {Статичний опис\\[-1pt]
   $H=\langle Z_v,Z_s,$\\
   $\rho_v,\rho_s,L_v,L_s,$\\
   $m_c,C_H\rangle$};
\node[dissertation small,text width=2.70cm] (behavior) at (-6.50,0.50)
  {Ознаки подій\\[-1pt]
   і поведінки $Z_b$};
\node[dissertation small,text width=3.10cm] (context) at (-6.50,-1.15)
  {Маска доступності,\\[-1pt]
   відомості й конфігурація\\
   $m_A,C_I,\theta_I$};

\node[dissertation small,text width=2.40cm,minimum height=19mm] (validation) at (-3.45,0.50)
  {Перевірка контракту\\
   ідентифікатор;\\
   контрольна сума;\\
   версії схем; походження;\\
   причинна\\
   належність;\\
   маска; якість};
\node[decision,minimum width=1.55cm,minimum height=9mm] (condition) at (-0.55,0.50)
  {Вимоги\\виконано?};
\node[dissertation small,text width=2.55cm,minimum height=11mm] (stop) at (-0.55,-1.75)
  {Інтеграцію припинено;\\
   причину відмови зафіксовано};
\node[dissertation small,text width=1.75cm,minimum height=10mm] (integration) at (2.10,0.50)
  {Узгоджене\\об’єднання\\$F_I$};
\node[decision,minimum width=1.45cm,minimum height=9mm] (composition) at (4.50,0.50)
  {Вхід\\повний?};
\node[dissertation small,text width=3.00cm,minimum height=11mm] (full) at (7.85,1.65)
  {Повне\\
   представлення\\[-1pt]
   $Z=\langle z,L,m_A,C_I\rangle$};
\node[dissertation small,text width=3.60cm,minimum height=13mm] (partial) at (7.85,-1.20)
  {Підтримуване часткове представлення\\[-1pt]
   $Z^{\partial}=\langle z_{sv},L_{sv},$\\
   $(m_s,m_v,0),C_I^{\partial}\rangle$};

\coordinate (inputbus) at (-4.80,0.50);
\draw[thick] (static.east) -- (static.east -| inputbus);
\draw[thick] (behavior.east) -- (behavior.east -| inputbus);
\draw[thick] (context.east) -- (context.east -| inputbus);
\draw[thick] (static.east -| inputbus) -- (context.east -| inputbus);
\fill[black] (inputbus) circle (1.3pt);
\draw[dissertation arrow] (inputbus) -- (validation.west);
\draw[dissertation arrow] (validation.east) -- (condition.west);
\draw[dissertation arrow] (condition.east) -- node[midway,above,font=\footnotesize,fill=white,inner sep=0.5pt]{так} (integration.west);
\draw[dissertation arrow] (condition.south) -- node[midway,right,font=\footnotesize,fill=white,inner sep=0.5pt]{ні} (stop.north);
\draw[dissertation arrow] (integration.east) -- (composition.west);
\coordinate (resultbus) at (5.75,0.50);
\draw[thick] (composition.east) -- (resultbus);
\draw[dissertation arrow] (resultbus) |- node[pos=0.20,right,font=\footnotesize,fill=white,inner sep=0.5pt]{так} (full.west);
\draw[dissertation arrow] (resultbus) |- node[pos=0.20,right,font=\footnotesize,fill=white,inner sep=0.5pt]{ні} (partial.west);
\end{tikzpicture}
\caption{Перевірка контракту та формування повного або часткового інтегрованого представлення}
\label{fig:feature_space}
\end{figure}


Отже, ознакове ядро утворює єдиний простір узгодження, у якому числові значення не відокремлюються від координат, доступності та походження. Завдяки цьому підтримувальне свідчення, суперечливе свідчення й відсутнє спостереження по-різному впливають на подальший висновок.

\section{Критерії прийняття рішення щодо наявності аномального коду}
\label{sec:decision_criteria}

Інтегроване представлення може містити сильні, але неповні свідчення або помірні, проте узгоджені спостереження. Зведення цих випадків до одного числа приховало б різну надійність підстав. Тому аналітичне рішення має окремо відображати клас стану, оцінку ризику, упевненість і пояснення.

Клас характеризує найближчий визначений тип стану. Ризик відображає критичність сукупності доступних свідчень. Упевненість показує повноту та стійкість підстав. Пояснення пов’язує результат зі спостережуваними ділянками й умовами перевірки. Захисна дія залежить від правил конкретної інформаційної технології, тому до виходу моделі не входить. Це розмежування дає змогу змінювати правила реагування без зміни вже одержаного аналітичного висновку.

Нехай $\Omega=\{K_0,K_1,\ldots,K_{n_A}\}$ --- множина класів стану, у якій $K_0$ відповідає недосягненню порога ознак аномального коду, а решта класів --- визначеним типам можливого прихованого впливу. Щоб усі складники результату належали одній перевірці та одному набору параметрів, їх формує єдина функція. Перехід від інтегрованого представлення до аналітичного стану наведено у формулі~\eqref{eq:analytic_state}:
\begin{equation}
S_A=F_D(Z,\tau),
\label{eq:analytic_state}
\end{equation}
де, $S_A$ --- аналітичний стан мультимедійного об’єкта; $F_D$ --- функція прийняття рішення; $Z$ --- інтегроване представлення; $\tau$ --- чинний поріг рішення.

Формула~\eqref{eq:analytic_state} установлює зв’язок між ознаками та завершеним результатом аналізу. До $S_A$ входять оцінки класів можливого впливу $p=(p_1,\ldots,p_{n_A})$, бінарний висновок $\hat y\in\{0,1\}$, визначений клас $\hat k$, інтегральна оцінка ризику $\rho_A$, упевненість $c_A$ і перевірюване пояснення $\Pi$. Класові оцінки не тлумачаться як статистичні ймовірності, доки їх калібрування не перевірено.

Часткові оцінки гілок обчислюються до інтеграції та не втрачаються під час об’єднання. Структурну оцінку $\rho_s$ і спектрально-зорову оцінку $\rho_v$ формує перший метод з ознак $Z_s$ і $Z_v$, подієво-поведінкову оцінку $\rho_b$ та оцінку прояву впливу $Q_A$ --- другий метод з ознак $Z_b$ і причинного контексту; правила їх обчислення встановлено в розділі~\ref{chap:methods_technology}. Усі чотири величини зберігаються серед іменованих складників вектора $z$, тому функція $F_D$ має до них безпосередній доступ. Для порівняння з порогом потрібна одна величина, яка зберігає внесок різних груп свідчень; її означено формулою~\eqref{eq:integral_risk}:
\begin{equation}
\rho_A=\omega_v\rho_v+\omega_s\rho_s+\omega_b\rho_b+\omega_q Q_A,
\label{eq:integral_risk}
\end{equation}
де, $\rho_v$, $\rho_s$ і $\rho_b$ --- часткові оцінки спектрально-зорової, структурної та подієво-поведінкової гілок; $Q_A$ --- оцінка прояву прихованого впливу через компонент веб-ресурсу; $\omega_v$, $\omega_s$, $\omega_b$ і $\omega_q$ --- вагові коефіцієнти; $\rho_A$ --- інтегральна оцінка ризику.

Часткові оцінки та $Q_A$ за побудовою набувають значень з відрізка $[0,1]$, а ваги є невід’ємними числами з одиничною сумою, тому й $\rho_A\in[0,1]$. Ваги відображають відносний внесок груп свідчень, а не ймовірність класу. Для недоступної гілки відповідну вагу встановлюють у нуль, після чого ваги доступних складників нормують за маскою $m_A$; оскільки оцінка $Q_A$ формується з подієво-поведінкової гілки, її доступність визначається складником $m_b$, і за $m_b=0$ вага $\omega_q$ обнуляється разом з $\omega_b$. Такий перерахунок дає змогу одержати частковий результат, але не робить його рівноцінним результату з повним набором спостережень: відмінність зменшує упевненість $c_A$ і відображається в поясненні.

Кандидатом класу є складник із найбільшою оцінкою, тобто $j^{*}=\arg\max_{j=1,\ldots,n_A}p_j$. Якщо $\rho_A<\tau$, призначається $K_0$ і $\hat y=0$; за умови $\rho_A\geq\tau$ призначається $K_{j^{*}}$ і $\hat y=1$. Однакові максимальні оцінки розв’язуються за наперед установленим порядком класів, який зберігається разом із налаштуваннями. Вектор $p$ характеризує лише розподіл між типами можливого впливу і не заміщує оцінки ризику: якщо $\rho_A\geq\tau$, а розподіл $p$ близький до рівномірного, визначений клас супроводжується низькою упевненістю, що обов’язково відображається в поясненні.

Бінарний висновок призначено для передавання результату до правила реагування, і його тлумачення обмежене. Значення $\hat y=1$ означає достатність доступних ознак для подальшої захисної дії, а $\hat y=0$ --- лише недосягнення чинного порога. Жодне зі значень не доводить виконання шкідливого коду й не гарантує безпечності.

Упевненість обчислюється як $c_A=F_q(Z,p)$ і залежить від близькості класових оцінок, доступності гілок за маскою $m_A$, якості даних, суперечливості ознак і причинної придатності подій, збережених у складниках $Z$. Низька упевненість не зменшує ризику автоматично. Високий ризик за низької упевненості означає наявність значущих, але недостатньо повних підстав; низький ризик за низької упевненості не є підтвердженням норми; високі значення обох показників також не замінюють перевірки програмного змісту.

Для відтворення причини рішення формується пояснення $\Pi=F_E(Z,\hat k,\rho_A,c_A)$. Структурне відхилення описується типом і байтовим проміжком, спектрально-зорове --- областю та часовою ділянкою, подієво-поведінкове --- фазою й причинним посиланням. Якщо ризик утворено поєднанням ознак, пояснення показує їх зв’язок, а не лише перелічує найбільші оцінки. До нього включаються також недоступні гілки, причини технічної непридатності, суперечності та версія порога.

До реалізацій функцій $F_D$, $F_q$ і $F_E$ висунуто три перевірювані вимоги. Вимога монотонності за доступністю: вилучення доступної гілки за незмінних інших умов не збільшує упевненості $c_A$. Вимога узгодженості перенормування: якщо всі доступні часткові оцінки рівні між собою, значення $\rho_A$ не залежить від складу доступних гілок. Вимога відтворюваності: за незмінних $X_{mm}$, версій засобів і параметрів повторне обчислення дає той самий аналітичний стан $S_A$. Жодна з вимог не потребує доведення в межах абстрактної моделі, але кожна перевіряється для конкретної реалізації, і порушення будь-якої з них позбавляє результат доказової сили.

Поділ аналітичного результату на клас, ризик, упевненість і пояснення відображено на рис.~\ref{fig:decision_criteria}. Він дає змогу перевіряти той самий $S_A$ незалежно від подальшого правила реагування.

\begin{figure}[htbp]
\centering
\begin{tikzpicture}[x=1cm,y=1cm]
\node[dissertation block,text width=2.85cm] (joint) at (-5.55,0.35)
  {Представлення\\$Z$};
\node[dissertation block,text width=2.85cm] (decision) at (-2.15,0.35)
  {Функція рішення\\$F_D$};
\node[dissertation block,text width=2.85cm] (state) at (1.25,0.35)
  {Аналітичний стан\\$S_A$};
\node[dissertation block,text width=3.40cm] (classification) at (5.05,2.15)
  {Визначення класу\\$p,\hat y,\hat k$};
\node[dissertation block,text width=3.40cm] (risk) at (5.05,0.35)
  {Ризик і упевненість\\$\rho_A,c_A$};
\node[dissertation block,text width=3.40cm] (explanation) at (5.05,-1.45)
  {Перевірюване пояснення\\$\Pi$};

\draw[dissertation arrow] (joint.east) -- (decision.west);
\draw[dissertation arrow] (decision.east) -- (state.west);
\draw[dissertation arrow] (state.east) -- ++(0.45,0) |- (classification.west);
\draw[dissertation arrow] (state.east) -- (risk.west);
\draw[dissertation arrow] (state.east) -- ++(0.45,0) |- (explanation.west);
\end{tikzpicture}
\caption{Формування аналітичного стану моделі без підміни його технологічною дією}
\label{fig:decision_criteria}
\end{figure}


Правило рішення забезпечує числову однозначність, але не приховує різниці між ризиком і надійністю підстав. Результат залишається обмеженим доступними даними та чинними налаштуваннями. Для переходу до розроблення методів залишається встановити порядок застосування моделі та умови передавання результатів між етапами.

\section{Процес застосування моделі до виявлення аномальних кодів}
\label{sec:detection_process_model}

Окремі частини моделі визначають дані, ознаки та рішення, однак без явного порядку їх реалізація могла б повторно змішати статичні й поведінкові дані або включити захисну дію до аналітичного стану. Завдання цього підрозділу полягає не у введенні ще однієї моделі, а у встановленні умов узгодження між моделлю, двома методами та інформаційною технологією розділу~\ref{chap:methods_technology}.

Порядок застосування побудовано за змістом виходів кожного етапу. Метод спектрально-зорового та структурного аналізу формує $Z_v$, $Z_s$, $L_v$ і $L_s$, після чого утворюється $H$. Метод поведінкового узгодження ознак різного походження формує $Z_b$ та виконує $F_I$. Функція $F_D$ завершує аналітичне оцінювання, а правило інформаційної технології окремо визначає захисну дію. Такий поділ пояснює призначення кожного етапу й не закріплює наперед конкретного способу програмної реалізації.

Межею завершення аналізу є перехід $S_A=F_D(Z,\tau)$, визначений формулою~\eqref{eq:analytic_state}. До його виконання мають бути підтверджені походження доступних представлень, сумісність схем ознак і маска доступності. Після виконання одержується аналітичний стан без захисної дії. Якщо передумову порушено, результат позначається як неповний або як технічна відмова, а не доповнюється даними іншої перевірки.

Щоб предметне тлумачення не змінювалося між переходами, встановлено шість незмінних правил:

\begin{enumerate}
\item \textit{Ідентичність входу.} Усі представлення однієї перевірки походять від того самого $X_{mm}$.
\item \textit{Незмінність.} Початковий мультимедійний об’єкт не заміщується похідним представленням.
\item \textit{Розділення даних.} Статичний опис не містить подієво-поведінкового представлення до етапу узгодженого об’єднання.
\item \textit{Явна доступність.} Відсутнє представлення не тлумачиться як нульова ознака.
\item \textit{Збереження походження.} Узагальнена ознака має зв’язок із байтовою, просторово-часовою або подієвою координатою.
\item \textit{Розмежування рішення і дії.} $S_A$ є аналітичним результатом, а захисна дія належить інформаційній технології.
\end{enumerate}

Ці правила визначають також граничні стани виконання. Кожний метод має завершуватися повним результатом, явно неповним результатом або відмовою. Неповний результат передається далі лише разом із маскою, показником якості та причиною обмеження. Відмова означає, що мінімальний вихід або його походження не підтверджено; вона не перетворюється на висновок про норму.

Модель не встановлює числових значень порогів, ваг або параметрів засобів кодування. Їх потрібно визначати на даних для налаштування та перевіряти на відокремлених даних. За відсутності такої перевірки можна описувати порядок обчислень, але не підтверджену перевагу або готовність до промислового застосування. Так само формально заповнений вихід не вважається коректним, якщо неможливо відтворити його походження чи версію схеми.

Для нового мультимедійного формату необхідно забезпечити незмінний $X_{mm}$, контрольовано одержані $X_v$, $X_f$, $X_s$, $X_m$ і, для відео, $X_t$, правильні значення $m_A$ та правила визначення координат. Ця вимога описує умову розширення моделі, але не підтверджує підтримки формату без реалізованого засобу аналізу й експериментальної перевірки. Якщо формат допускає кілька коректних тлумачень, у $C$ зберігається використане правило, а у висновку --- межа його застосовності.

Межу відповідальності між моделлю, двома методами та п’ятиетапною інформаційною технологією показано на рис.~\ref{fig:detection_process}. Таке розмежування завершує модельний рівень і дає змогу конкретизувати обчислення та реагування без уведення додаткової моделі або зміни змісту $S_A$.

\begin{figure}[htbp]
\centering
\begin{tikzpicture}[x=1cm,y=1cm]
\node[dissertation block,text width=9.8cm,minimum height=12mm] (model) at (0,3.6)
  {Модель мультимодального представлення\\прихованої загрози};

\node[dissertation block,text width=5.2cm,minimum height=13mm] (method1) at (-3.6,1.2)
  {Метод аналізу\\спектральних, зорових\\і структурних ознак};
\node[dissertation block,text width=5.2cm,minimum height=13mm] (method2) at (3.6,1.2)
  {Метод поведінкового\\узгодження ознак\\різного походження};

\node[dissertation block,text width=9.8cm,minimum height=12mm] (technology) at (0,-1.4)
  {П’ятиетапна інформаційна технологія\\виявлення ознак аномального вмісту};

\draw[dissertation arrow] (model.south) -- ++(0,-0.35) -| (method1.north);
\draw[dissertation arrow] (model.south) -- ++(0,-0.35) -| (method2.north);
\draw[dissertation arrow] (method1.east) -- (method2.west);
\draw[dissertation arrow] (method1.south) -- ++(0,-0.35) -| (technology.north);
\draw[dissertation arrow] (method2.south) -- ++(0,-0.35) -| (technology.north);
\end{tikzpicture}
\caption{Взаємозв’язок моделі, двох методів та інформаційної технології}
\label{fig:detection_process}
\end{figure}


Процес застосування не дублює внутрішніх процедур методів: він визначає послідовність переходів, межі довіри й граничні стани, за яких результати можна узгоджувати та передавати до правила реагування. У розділі~\ref{chap:methods_technology} ці умови конкретизовано двома методами та п’ятиетапною інформаційною технологією без зміни змісту модельних виходів.

\section*{Висновки до розділу 2}
\addcontentsline{toc}{section}{Висновки до розділу 2}

У розділі запропоновано розв’язання дослідницької суперечності між потребою врахувати слабкі різнорідні ознаки та необхідністю обмежити силу висновку фактично доступними свідченнями: розроблено модель мультимодального представлення прихованої загрози в мультимедійних об’єктах, у якій умови здійснення атаки, багаторівневе представлення, ознакове ядро, критерії рішення та процес застосування є взаємозалежними частинами одного цілого.

Предметне тлумачення результату ґрунтується на розмежуванні аномалії, аномального коду, прихованої загрози та шкідливого коду, а також супровідних понять прихованого навантаження і прихованого впливу. Завдяки цьому форматне або статистичне відхилення не ототожнюється з підтвердженою несанкціонованою дією.

Початковий і підготовлений мультимедійні об’єкти $X_{mm}$ та $X_{mm}^{\prime}$ пов’язано з даними $X_v$, $X_f$, $X_s$, $X_m$, $X_t$, $X_b$, ознаками $Z_v$, $Z_s$, $Z_b$ та інтегрованим представленням $Z$. Маска $m_A$ відокремлює відсутнє спостереження від виміряного нульового значення, а координати й контекст зберігають якість і походження ознак. Межу довіри проведено між незмінним $X_{mm}$ та результатами його тлумачення із зазначенням використаних засобів і їх версій.

Аналітичний стан $S_A$ містить класові оцінки, бінарний висновок, визначений клас, оцінку ризику, упевненість і пояснення. Ризик відокремлено від упевненості, тому неповнота або суперечність підстав не приховується одним пороговим значенням. Захисну дію вилучено з модельного стану, оскільки її визначає інформаційна технологія; це розмежування утворює однозначний перехід до розділу~\ref{chap:methods_technology}.

Межами одержаного результату є залежність від підтримуваних форматів, якості контрольованого декодування, повноти причинно пов’язаних подій та перевіреності параметрів. Модель не гарантує виявлення довільного прихованого навантаження й не підтверджує промислової готовності без експериментальної перевірки методів та інформаційної технології.
